\documentclass[12pt]{article}
\usepackage[a4paper, tmargin=3cm, lmargin=2.5cm, rmargin=2.5cm, bmargin=3cm]{geometry}
\usepackage[utf8]{inputenc}
\usepackage[T1]{fontenc}
\usepackage{hyperref}

\usepackage[brazilian]{babel} % Remover linha em caso de escrita em inglês
\usepackage[fixlanguage]{babelbib} % Remover linha em caso de escrita em inglês


\hypersetup{
    colorlinks=true,
    linkcolor=black,
    filecolor=magenta,      
    urlcolor=blue,
    citecolor=black,
}

\usepackage{natbib} % to cite references by surname and year
\usepackage{graphicx}
\usepackage{fancyhdr}
\pagestyle{fancy} % Ativar o estilo fancy
\fancyhf{} % Limpar cabeçalhos e rodapés
\renewcommand{\headrulewidth}{0pt} % Remover a linha do cabeçalho
\fancyfoot[C]{IC/UNICAMP \\ 2026 }


%%%%%%%%%%%%%%%%%%%%%%%%%%%%%%%%%%
% CONFIGURAÇÃO DA PRIMEIRA PÁGINA
%%%%%%%%%%%%%%%%%%%%%%%%%%%%%%%%%%


\begin{document}

\vspace{5mm}
\hrule

\vspace{3mm}

\begin{center}
    {\large\textbf{INSTITUTO DE COMPUTAÇÃO - UNICAMP}} \\
    \vspace{3mm}
    {\large\textbf{MO810A/MC959A - Sistemas Multiagentes com Modelos de Linguagem}} \\
    \vspace{3mm}
    {\large\textbf{Prof. Dr. Julio Cesar dos Reis}}
\end{center}

\vspace{3mm}
\hrule

\vspace{2.5cm}
\begin{flushright}
    \Large{\textbf{Fase 1 - Proposta}} \\
    \vspace{1cm}
    \textit{
        Does Topology Matter? Auditing Socioeconomic Bias Across Multi-Agent LLM Architectures in Hiring and Criminal Verdicts \\
    }
    \vspace{5cm}
    Luiz Felipe Lenharo - 237896 \\
    Felipe Rocha Verol - 248552 
\end{flushright}

%%%%%%%%%%%%%%%%%%%%%%%%%%%%%%%%%%
%%%%%%%%%%%%%%%%%%%%%%%%%%%%%%%%%%


\newpage
\pagestyle{fancy} % 2. Ativar o estilo fancy
\fancyhf{} % 3. Limpar cabeçalhos e rodapés
\renewcommand{\headrulewidth}{0pt} % Remover a linha do cabeçalho

\fancyfoot[R]{\thepage}


%%%%%%%%%%%%%%%%%%%%%%%%%%%%%%%%%%
%       INÍCIO DAS SEÇÕES
%%%%%%%%%%%%%%%%%%%%%%%%%%%%%%%%%%

%%%%%%%%%%%%%%%%%%%%%%%%%%%%%%%%%%%%%%%%
\section{Cenário e Contexto}
\label{sec:contexto}
%%%%%%%%%%%%%%%%%%%%%%%%%%%%%%%%%%%%%%%%

Este trabalho se insere na área de IA aplicada à justiça social e equidade algorítmica, especificamente na auditoria de viés em sistemas multiagentes baseados em Modelos de Linguagem (LLMs). O tema se conecta diretamente ao ODS 10 (Redução das Desigualdades), com aplicação em dois domínios que instanciam outros dois ODS: avaliação de candidatos a vagas de emprego (ODS 8 -- Trabalho Decente e Crescimento Econômico) e apoio à decisão em veredictos de júri penal (ODS 16 -- Paz, Justiça e Instituições Eficazes).

O ambiente em que o problema ocorre é o de organizações -- empresas de recrutamento, plataformas de triagem de currículos, ferramentas de apoio a decisões jurídicas -- que vêm delegando etapas de pré-análise de casos sensíveis a sistemas baseados em LLM. Diferente de uma única chamada a um modelo, esses sistemas são cada vez mais construídos como arquiteturas multiagentes, em que múltiplos agentes (por exemplo, extrator de perfil, avaliador, crítico/revisor, agregador de opiniões) cooperam ou competem para produzir uma decisão final. Frameworks como AutoGen e MetaGPT popularizaram esse tipo de orquestração, e há evidências de que a topologia de orquestração (pipeline, hierárquica, debate, ensemble, etc.) afeta métricas como custo, robustez e taxa de falha da tarefa \textbf{[FALTA: referência específica -- ex. artigo "The Harness Effect" da lista Awesome Agentic Systems discute custo/token economics, mas precisa confirmar se é aplicável aqui]}.

Paralelamente, existe uma linha de pesquisa consolidada demonstrando que LLMs isolados replicam vieses socioeconômicos, raciais e de gênero quando avaliam currículos ou casos jurídicos \textbf{[FALTA: citações específicas -- ex. estudos de viés em triagem de currículos por LLM e estudos de viés racial/socioeconômico em sistemas de apoio a decisão penal (tipo COMPA/COMPAS na literatura de ML tradicional). Não inventei essas referências porque não tenho certeza dos títulos/autores/anos exatos -- precisa localizar e confirmar antes de citar]}. O que ainda não está bem estabelecido na literatura é como esse viés se comporta quando o mesmo modelo-base é inserido em diferentes arquiteturas multiagentes: se a orquestração atenua, amplifica ou não altera o viés observado em avaliações de agente único.

%%%%%%%%%%%%%%%%%%%%%%%%%%%%%%%%%%%%%%%%
\section{Motivação para o Tema}
\label{sec:motivacao}
%%%%%%%%%%%%%%%%%%%%%%%%%%%%%%%%%%%%%%%%

A escolha do tema é motivada por três fatores. Primeiro, um fator social: a desigualdade socioeconômica é um problema crescente e mensurável, e decisões automatizadas em contratação e justiça penal afetam diretamente o acesso de grupos vulneráveis a oportunidades econômicas e a tratamento justo perante a lei -- conexão direta com o ODS 10. Segundo, um fator científico/tecnológico: há preocupação documentada de que LLMs reproduzam vieses presentes nos dados de treinamento e no comportamento de usuários \textbf{[FALTA: referência]}, mas a maior parte da auditoria de viés em LLM foca em modelo único; falta entender se arquiteturas multiagentes (cada vez mais adotadas em produtos reais de HR-tech e legal-tech) mudam esse comportamento. Terceiro, um fator prático: entender esse comportamento é pré-requisito para responder, com rigor, se sistemas multiagentes são de fato adequados para pré-análise de casos sensíveis -- hoje isso é uma pergunta em aberto, não um pressuposto do projeto.

O impacto esperado tem dois grupos de beneficiários e mecanismos distintos: (i) candidatos a emprego e réus/acusados, que seriam protegidos indiretamente caso o benchmark leve à identificação e ao desencorajamento de arquiteturas que amplificam viés discriminatório; e (ii) equipes de engenharia e reguladores que precisam escolher ou auditar uma topologia de orquestração antes de colocá-la em produção, para os quais o benchmark funcionaria como uma metodologia de comparação replicável entre arquiteturas.


%%%%%%%%%%%%%%%%%%%%%%%%%%%%%%%%%%%%%%%%
\section{Caracterização do Problema}
\label{sec:problema}
%%%%%%%%%%%%%%%%%%%%%%%%%%%%%%%%%%%%%%%%

O problema investigado é: não se sabe se, e em que grau, o viés socioeconômico observado em avaliações feitas por um único agente LLM é atenuado, preservado ou amplificado quando o mesmo modelo-base é organizado em diferentes topologias de sistemas multiagentes, nas tarefas de triagem de candidatos a emprego e de recomendação em veredictos de júri penal.

A lacuna se sustenta em duas frentes de evidência que, isoladamente, não respondem à pergunta: de um lado, há estudos que documentam viés em avaliações de agente único por LLM \textbf{[FALTA: referências específicas de viés em triagem de currículo e/ou decisão penal por LLM]}; de outro, há estudos que mostram que a topologia de orquestração multiagente altera métricas de desempenho, custo e robustez da tarefa (ex.: AutoGen, MetaGPT, e trabalhos sobre economia de tokens em diferentes arquiteturas de orquestração \textbf{[FALTA: confirmar quais desses artigos foram efetivamente lidos]}), mas nenhum desses trabalhos mede viés socioeconômico como variável de interesse. A interseção entre as duas linhas -- efeito da topologia sobre o viés, especificamente -- \textbf{[FALTA: confirmar se essa lacuna é explicitamente apontada em algum survey/related work, ou se a alegação de ineditismo se sustenta apenas por ausência de resultado nas buscas realizadas]}.

As consequências de não resolver essa lacuna são práticas e científicas: praticamente, produtos de HR-tech e legal-tech que adotam múltiplos agentes podem estar transmitindo uma falsa sensação de segurança (``mais agentes deliberando = decisão mais justa'') sem evidência que sustente essa suposição; cientificamente, falta uma metodologia comparável para avaliar como escolhas de design de orquestração interagem com equidade algorítmica.

Escopo -- dentro do projeto: (i) avaliação controlada via comparação pareada contrafactual, variando uma única característica correlacionada a status socioeconômico por vez e mantendo as demais constantes; (ii) duas tarefas de decisão (triagem de candidatos, recomendação de veredicto penal); (iii) um conjunto definido de arquiteturas multiagentes, incluindo uma linha de base de agente único como referência de comparação. Fora do escopo: (i) implantação ou teste com usuários reais; (ii) fine-tuning ou alteração dos modelos de linguagem subjacentes; (iii) explicação causal do mecanismo interno pelo qual a topologia influencia o viés (o projeto mede \textit{se} e \textit{quanto}, não necessariamente \textit{por quê}); (iv) eixos de viés não relacionados a status socioeconômico, exceto quando necessários para caracterizar interseccionalidade -- \textbf{[FALTA: decisão do grupo sobre incluir ou não raça/gênero como eixos adicionais]}.

%%%%%%%%%%%%%%%%%%%%%%%%%%%%%%%%%%%%%%%%
\section{Proposta de Solução Conceitual}
\label{sec:solucao}
%%%%%%%%%%%%%%%%%%%%%%%%%%%%%%%%%%%%%%%%

A proposta é um benchmark de auditoria de viés por comparação pareada contrafactual, aplicado a um conjunto de arquiteturas multiagentes candidatas. Para cada tarefa (triagem de emprego e veredito penal), o benchmark gera pares de casos idênticos exceto por uma única característica correlacionada a status socioeconômico (ex.: bairro/CEP como proxy de renda, nome do candidato, prestígio da instituição de ensino, histórico penal contextualizado). Cada arquitetura avalia todos os casos e atribui uma pontuação ou decisão; o benchmark compara, par a par, a diferença de pontuação atribuída a casos que só diferem naquela característica.

As arquiteturas candidatas -- cada uma contribuindo para isolar um mecanismo diferente pelo qual o viés poderia se propagar -- são \textbf{[FALTA: definição final do grupo]}, mas incluem candidatas como:
\begin{itemize}
    \item \textbf{Agente único} (linha de base): um único agente LLM realiza a avaliação completa, servindo como referência para medir o efeito líquido de introduzir múltiplos agentes.
    \item \textbf{Pipeline sequencial}: agentes especializados (ex.: extração de perfil, avaliação, síntese) atuam em série, testando se a decomposição em etapas concentra ou dilui viés presente em uma etapa específica.
    \item \textbf{Debate/adversarial}: dois agentes argumentam posições opostas e um terceiro agente julga, testando se o confronto de perspectivas expõe ou neutraliza um viés individual.
    \item \textbf{Revisão hierárquica}: um agente avaliador é supervisionado por um agente revisor com poder de veto/ajuste, testando se a supervisão corrige ou herda o viés do avaliador.
    \item \textbf{Ensemble/votação}: múltiplos agentes avaliam independentemente e o resultado é agregado (média ou voto), testando se a agregação estatística atenua vieses individuais.
\end{itemize}

A contribuição de cada arquitetura para a solução é, portanto, permitir isolar um mecanismo estrutural específico (especialização em série, confronto adversarial, supervisão hierárquica, agregação estatística) e compará-lo contra a linha de base de agente único, sem a qual não é possível atribuir uma diferença observada à topologia em si. A métrica de viés usada para comparar os pares -- \textbf{[FALTA: definição formal, ex.: diferença média de pontuação, razão de impacto adverso, teste de significância estatística considerando múltiplas execuções por caso dada a estocasticidade do LLM]} -- é aplicada de forma idêntica a todas as arquiteturas, permitindo comparação direta entre elas.


%%%%%%%%%%%%%%%%%%%%%%%%%%%%%%%%%%%%%%%%
\section{Posicionamento em Relação ao Estado da Arte}
\label{sec:estado_da_arte}
%%%%%%%%%%%%%%%%%%%%%%%%%%%%%%%%%%%%%%%%

\textbf{[SEÇÃO INCOMPLETA -- NÃO PREENCHER COM CITAÇÕES INVENTADAS. Ver lista de pendências enviada junto com este rascunho.]}

Da lista \href{https://juliodosreis.com/awesome-agenticsystems/}{Awesome Agentic Systems}, nenhum artigo trata de viés/equidade -- a lista é focada em arquitetura, memória, ferramentas e avaliação de agentes. Os candidatos mais próximos do projeto, por tratarem de topologia de orquestração multiagente (relevantes à Seção~\ref{sec:solucao}, não à questão de viés em si), são:
\begin{itemize}
    \item \textbf{AutoGen: Enabling Next-Gen LLM Applications via Multi-Agent Conversation} (Wu et al., 2023) -- modela a aplicação como conversa entre agentes configuráveis; candidato a referência para justificar a arquitetura de debate/pipeline.
    \item \textbf{MetaGPT: Meta Programming for a Multi-Agent Collaborative Framework} (Hong et al., 2023) -- codifica papéis de agente conectados por artefatos estruturados; candidato para justificar a arquitetura de pipeline/hierárquica.
    \item \textbf{The Harness Effect: How Orchestration Design Sets the Token Economics of Enterprise Agentic AI} (Sayed Ali et al., 2026) -- mostra que trocar a camada de orquestração, mantendo o modelo fixo, muda custo/robustez; evidência indireta de que topologia importa, mas não trata de viés.
\end{itemize}

\textbf{[FALTA: confirmar quais desses três (ou outros da lista) foram de fato lidos pelo grupo -- o template exige citar apenas artigos lidos.]} Para cada artigo efetivamente selecionado, falta ainda preencher, por artigo: (i) critério de seleção; (ii) o que foi extraído; (iii) que adaptação o método exige para tratar o problema da Seção~\ref{sec:problema}.

Além disso, nenhum artigo da lista trata de viés algorítmico, equidade em triagem de currículos ou justiça penal computacional -- essas referências precisam vir de fora da lista (conforme permitido pelo enunciado) e ainda não foram levantadas neste rascunho.

%%%%%%%%%%%%%%%%%%%%%%%%%%%%%%%%%%%%%%%%
\section*{Apêndice}
\label{sec:apendice}
%%%%%%%%%%%%%%%%%%%%%%%%%%%%%%%%%%%%%%%%

\textit{$<$Use se necessário. $>$}



%%%%%%%%%%%%%%%%%%%%%%%%%%%%%%%%%%
%          BIBLIOGRAFIA
%%%%%%%%%%%%%%%%%%%%%%%%%%%%%%%%%%
\newpage
\bibliographystyle{aasjournal}
\bibliography{references}{}



\end{document}
